\section{Preliminaries}\label{sec:preliminaries}

\subsection{Tensor products and positive elements}
All algebras are complex. All tensor products in the construction are
\emph{maximal}, and we abbreviate $D^{\otimes_{\max}n}$ to $D^{\otimes n}$.
We specify minimal tensor products explicitly. Algebraic tensor products
are denoted by $\odot$ over $\C$ and by $\otimes_{\R}$ over $\R$.
Representations called irreducible are nonzero and nondegenerate, hence
unital when the represented algebra is unital. An element is \emph{full}
if its generated closed two-sided ideal is the whole algebra. Equivalently,
it is nonzero in every irreducible representation: a proper ideal has a
nonzero irreducible representation of its quotient.

For positive elements in finite matrix algebras over a unital algebra $B$,
write $a\precsim b$ if, after adjoining zero corners as necessary, there
are matrices $x_j$ with $x_j^*b x_j\to a$ in norm. Write $[a]$ for the
resulting Cuntz equivalence class. Thus $[a]\le[b]$ means
$a\precsim b$, and addition is induced by direct sum. We use only the
following elementary rules:
\begin{enumerate}
\item $[v^*v]=[vv^*]$ for every rectangular matrix $v$;
\item if $a_1,\ldots,a_l$ are positive and pairwise orthogonal, then
$[\sum_i a_i]=\sum_i[a_i]$;
\item comparison and addition are preserved by $*$-homomorphisms and by
external tensoring with a fixed positive element, including for maximal
tensor products.
\end{enumerate}
For the first rule, with $b=v^*v$, the matrices
$x_\gamma=v(b+\gamma)^{-1/2}$ satisfy
$x_\gamma^*(vv^*)x_\gamma=b^2(b+\gamma)^{-1}\to b$ as
$\gamma\downarrow0$; apply the same argument to $v^*$ for the converse.
For the second rule, apply the first to the row
$(a_1^{1/2},\ldots,a_l^{1/2})$. For the third rule, apply a homomorphism to
the witnesses, or tensor them with the unit of the other algebra. The norm
estimate for the resulting error is at most the original error times the
norm of the fixed positive tensor factor. Direct sums agree under the
canonical matrix identifications. These rules use no cancellation or
regularity of the Cuntz order. We also use $[a]\le[1_B]$ for $a\in B_+$,
which follows from the exact witness $a^{1/2}$.

\subsection{The dimension-drop criterion}
We say that positive elements $a_1,a_2\in B$ are \emph{equivalent} if
$a_1=v^*v$ and $a_2=vv^*$ for one $v\in B$. This is the exact relation
required in the next proposition.

\begin{proposition}[R\o rdam--Winter]\label{prop:rw}
Let $B$ be a unital $C^*$-algebra. Suppose that $a_1,a_2\in B$ are
orthogonal positive contractions with $a_1=v^*v$, $a_2=vv^*$ for some
$v\in B$, and that, for some $\varepsilon>0$,
\[
 1_B-a_1-a_2\precsim(a_1-\varepsilon)_+.
\]
Then there is a unital $*$-homomorphism $I(2,3)\to B$.
\end{proposition}

This is \cite[Proposition~5.1, (ii)$\Rightarrow$(iv), with $n=2$]{RW}.
Section~\ref{sec:rw} supplies a complete proof using their method, with
the support projections, scalar endpoint fiber, and weighted flip
identity made explicit.
The stable-rank-one hypothesis appearing elsewhere in that proposition is
needed only for (i)$\Rightarrow$(ii). It is not a hypothesis of the
implication stated here. Orthogonality ensures $a_1+a_2\le1_B$, so the
defect in the display is positive.

\subsection{Central-sequence inputs}
We use the following results of Kirchberg and R\o rdam for the fixed free
ultrafilter $\omega$.

\begin{proposition}[Kirchberg--R\o rdam]\label{prop:kr}
Let $A$ be a nonzero unital separable $C^*$-algebra.
\begin{enumerate}
\item If $D\subset F_\omega(A)$ is separable and contains the ambient
unit, then for every $m\ge1$ there is a unital $*$-homomorphism
$D^{\otimes_{\max}m}\to F_\omega(A)$.
\item $A\cong A\otmin\Z$ if and only if $F_\omega(A)$ contains, with
the ambient unit, a separable subhomogeneous $C^*$-algebra without
characters.
\end{enumerate}
\end{proposition}

The first assertion is the unital specialization of the commuting-copies
lemma in \cite[Lemma~9.3]{Fixed231}, whose primary source is
\cite[Proposition~2.2 and Corollary~2.3]{KR}. For unital $A$, the
annihilator ideal used in that companion is zero: an element annihilating
every element of $A$ annihilates its unit. Thus its $F(A)$ is precisely
the norm central-sequence algebra used here. This particular lemma has
no nuclear-dimension assumption. We include the short specialization
to make its quantifiers explicit.

\begin{proof}[Proof of Proposition~\ref{prop:kr}(1)]
For any separable subalgebra $E\subset F_\omega(A)$, we construct a
unital map from $D$ into $F_\omega(A)\cap E'$. Choose a countable dense
unital $*$-subalgebra $D_0\subset D$ over $\mathbb Q+i\mathbb Q$.
For every $x\in D_0$, choose a bounded lift $(a_l^x)_l$ in $A$ with
$\sup_l\|a_l^x\|\le2\|x\|$, using the zero lift for $x=0$ and the
constant unit lift for $x=1$. Such choices follow from the quotient norm,
by truncating coordinate norms without changing the ultrapower class.
All additive, rational-complex scalar, multiplicative and adjoint
relation defects tend to zero in norm along $\omega$.

Choose a countable dense subset of $A$ and bounded lifts $(b_s^j)_s$
of a countable dense subset of $E$. Enumerate the relation-defect tests
and the tests for centrality of the $a_l^x$ against the chosen subset
of $A$. For each output index $s$, choose $l(s)$ so that the first $s$
tests hold within $1/s$, and so that
\[
 \|[a_{l(s)}^x,b_s^j]\|<1/s
\]
for the first $s$ pairs $(x,j)$ in an enumeration. Each of the finitely
many requirements holds on an $\omega$-large set of old indices $l$:
for this fixed $s$, each $b_s^j$ is a fixed element of $A$, so original
centrality applies. Their finite intersection is nonempty.

The new sequences $(a_{l(s)}^x)_s$ are central and satisfy every
algebraic relation in the ordinary limit, hence along the arbitrary
free ultrafilter $\omega$. Their classes define a unital
rational-complex $*$-homomorphism on $D_0$ into $F_\omega(A)\cap E'$.
The uniform bound $2\|x\|$ extends it continuously to a complex unital
$*$-homomorphism on $D$. Starting with the given copy of $D$ and taking
$E$ successively to be the algebra generated by the earlier images
produces any prescribed finite number of commuting unital images.
The universal property of the maximal tensor product gives (1).
Injectivity of the reindexed maps is not needed.
\end{proof}

The second assertion is \cite[Theorem~2.5]{KR}. Here
\emph{subhomogeneous} means that the dimensions of all irreducible
representations have a common finite bound. In particular, this input
does not replace norm central sequences by a tracial ultrapower or
impose nuclearity on $A$.
